\section{Alternative proof of the single-layer bound via Milnor--Thom}\label{app:milnor-thom}

The proof of Proposition~\ref{prop:single-layer-degree} given in Section~\ref{sec:tubular-neural} uses Bernstein's theorem to obtain the sharp bound $\md(V) = O(w^n)$. In this section we present a weaker alterantive argument based on Milnor--Thom dimension-climbing. While it only yields a weaker bound $\md(V) = O(w^{n^2})$, the argument extends more readily to settings in which the system does not admit a common Newton polytope (e.g., Pfaffian systems that are not polynomial after an exponential substitution).

The proof hinges on a per-line bound on the zeros of a sum of sigmoid derivatives (Lemma~\ref{lem:line-intersection}), based on the P\'olya--Voorhoeve theorem for exponential polynomials (Lemma~\ref{lem:polya}). While this result is not used in the BKK proof of Proposition~\ref{prop:single-layer-degree}, it is of independent interest and serves as main ingredient of an alternative, weaker proof via Milnor--Thom dimension-climbing that we defer to Appendix~\ref{app:milnor-thom}.

\begin{lemma}\label{lem:polya}
    A non-trivial exponential polynomial $\varphi(t) = \sum_{i=1}^M d_i \econst^{\nu_i t}$ with $M$ pairwise distinct real frequencies has at most $M-1$ real zeros.
\end{lemma}

\begin{proof}
The proof is by induction on $M$. For $M=1$, the function $\varphi(t)$ has $M-1=0$ zeros. Assume the result holds for $M-1$. The function $\mathrm{e}^{-\nu_1 t}\varphi(t)$ has the same number of zeros as $\varphi(t)$ and its derivative is the sum of $M-1$ terms, so by the induction hypothesis has at most $M-2$ zeros. The claim now follows by Rolle's theorem.
\end{proof}

\begin{lemma}\label{lem:line-intersection}
Let $\sigma(t) = (1+\econst^{-t})^{-1}$ be the logistic sigmoid, and let $a_1,\ldots,a_w\in \mathbb{Q}^n$ with common denominator $q\in \Z_{>0}$ and set $L := q\cdot \max_{k,i}|a_{ki}|$. Let $b_1,\ldots,b_w,\gamma_1,\ldots,\gamma_w,c\in \R$ be constants, and consider the function
\begin{equation*}
  h(x) = c + \sum_{k=1}^w \gamma_k \sigma'(a_k^{\trans}x + b_k).
\end{equation*}
Then for any line $\ell\subset \R^n$ in general position, the restriction $h|_\ell$ has at most
\begin{equation*}
  (4Lw+1)^n - 1
\end{equation*}
non-degenerate real zeros.
\end{lemma}

\begin{remark}\label{rem:commensurability-sufficient}
    The rationality hypothesis $a_k\in \mathbb{Q}^n$ in Lemma~\ref{lem:line-intersection} is automatic for any finite-precision (e.g.,\ floating-point) weight matrix: taking $q$ to be the least common multiple of the denominators gives $q a_{ki}\in\Z$ with $L = q\,\max_{k,i}|a_{ki}|$ bounded by the machine precision.
    \end{remark}

\begin{proof}[Proof of Lemma~\ref{lem:line-intersection}]
Let $\ell(t) = x_0 + te$ parametrise the line. Set $\lambda_k = a_k^{\trans} e$, $\mu_k = a_k^{\trans} x_0 + b_k$, and $g(t) := h(\ell(t))$. Since $\sigma'(s) = \econst^{-s}/(1+\econst^{-s})^2$, the substitution $u_k(t) = \econst^{-\lambda_k t - \mu_k}$ gives
\begin{equation*}
  \sigma'(\lambda_k t + \mu_k) = \frac{u_k}{(1+u_k)^2}.
\end{equation*}
Multiplying $g(t)$ by the positive factor $\Phi(t) = \prod_{k=1}^w (1+u_k(t))^2$ preserves zeros and yields
\begin{equation}\label{eq:Htilde}
  \widetilde{g}(t) := g(t)\cdot \Phi(t) = c\prod_{k=1}^w (1+u_k)^2 + \sum_{k=1}^w \gamma_k\, u_k\prod_{j\neq k}(1+u_j)^2.
\end{equation}
Expanding in terms of monomials $u_1^{s_1}\cdots u_w^{s_w}$ for $s_k\in \{0,1,2\}$, and using $u_k(t) = \econst^{-\lambda_k t - \mu_k}$, we rewrite~\eqref{eq:Htilde} as the exponential polynomial
\begin{equation}\label{eq:gen-exp-poly}
  \widetilde{g}(t) = \sum_{s \in \Lambda} d_{s}\, \econst^{-\nu_{s} t},
\end{equation}
where $s$ ranges over those multi-indices in $\{0,1,2\}^w$ with nonzero combined coefficient, $\nu_s = \sum_{j=1}^w s_j \lambda_j$, and $d_s$ collects the contributions of all monomials sharing the frequency $\nu_s$. Lemma~\ref{lem:polya} bounds the number of zeros by $|\Lambda| - 1$, so it remains to bound $|\Lambda|$, the number of distinct frequencies.

Let $V := \mathrm{span}_{\mathbb{Q}}(\lambda_1,\ldots,\lambda_w) \subset \R$ and assume that $e$ is such that $\dim \mathrm{span}_{\mathbb{Q}}(e_1,\ldots,e_n) = n$ (this is the genericity assumption in the statement, as it holds for almost all $e$). Since each $\lambda_k = \sum_{i=1}^n a_{ki} e_i$ is a $\mathbb{Q}$-linear combination of $e_1,\ldots,e_n$, we have $\dim_\mathbb{Q} V \leq n$. Taking the $\mathbb{Q}$-basis $\lambda_j' := e_j/q$ of the ambient lattice $\mathbb{Z}\lambda_1'\oplus\cdots\oplus\mathbb{Z}\lambda_n'$, the identity $\lambda_k = \sum_i a_{ki}e_i = \sum_i (q a_{ki})\,\lambda_i'$ expresses each $\lambda_k$ as an integer combination of the $\lambda_j'$ with coefficients $m_{kj} := q a_{kj}\in \Z$ satisfying $|m_{kj}|\leq L$. Consequently
\begin{equation*}
  \nu_{s} \;=\; \sum_{j=1}^n \Big(\sum_{k=1}^w s_k m_{kj}\Big)\lambda_j',
\end{equation*}
and since $\lambda_1',\ldots,\lambda_n'$ are $\mathbb{Q}$-linearly independent, $\nu_{s} = \nu_{s'}$ if and only if $\sum_k s_k m_{kj} = \sum_k s_k' m_{kj}$ for all $j\in [n]$. Hence the number of distinct frequencies equals the cardinality of the image of
\begin{equation}\label{eq:freq-lattice-map}
  \vec n\colon \{0,1,2\}^w \to \Z^n, \qquad \vec n(s) \;=\; \Big(\sum_{k=1}^w s_k m_{k1},\; \ldots,\; \sum_{k=1}^w s_k m_{kn}\Big).
\end{equation}
Each coordinate of $\vec n(s)$ lies in $[-2Lw, 2Lw]$, so
\begin{equation*}
  |\Lambda| \;\leq\; (4Lw+1)^n,
\end{equation*}
and Lemma~\ref{lem:polya} gives $Z(h|_\ell) = Z(g) = Z(\widetilde{g}) \leq (4Lw+1)^n - 1$.
\end{proof}

% \begin{remark}[Zonotope interpretation]\label{rem:zonotope}
% The image of the map~\eqref{eq:freq-lattice-map} counts lattice points in the Minkowski sum
% \begin{equation*}
%   \sum_{k=1}^w \bigl\{\,0,\, \vec m_k,\, 2\vec m_k\,\bigr\}, \qquad \vec m_k = (m_{k1},\ldots,m_{kn}) \in \Z^n,
% \end{equation*}
% which is contained in the zonotope $Z = \sum_{k=1}^w [0,2\vec m_k]\subset \R^n$. For coefficients $\vec m_k$ in general position with $\|\vec m_k\|_\infty \leq L$, the number of lattice points in $Z$ is $\Theta(w^n)$, with a constant depending on the $\vec m_k$. The bound $|\Lambda|\leq (4Lw+1)^n$ is therefore tight up to constants, and the exponent $n$ is sharp. This matches the lower bound $\Omega(w^n)$ for $\md(V_{ij})$ established in Remark~\ref{rem:single-layer-lowerbound}. It is plausible that a refinement of the zero-counting bound via mixed volumes and the BKK theorem (applied to exponential sums with commensurable frequencies supported in $\{0,1,2\}^w$) could close the gap between the upper bound $O(w^{n^2})$ in Proposition~\ref{prop:single-layer-degree} and the lower bound $\Omega(w^n)$, but we do not pursue this here.
% \end{remark}

% \begin{remark}\label{rem:line-intersection-pfaffian}
% In geometric terms, the zeros of $h(\ell(t))$ correspond to intersections of the curve $\Gamma = \{(s_1(t),\ldots,s_w(t))\}\subset (0,1)^w$, $s_k(t) = \sigma(\lambda_k t+\mu_k)$, with the hypersurface $\{c+\sum_k \gamma_k P(s_k)=0\}$. This is an instance of the Khovanski\u i--Rolle framework (Lemma~\ref{lem:kho-rolle}), with $\Gamma$ playing the role of the curved $1$-manifold. Applying the standard Khovanski\u i bound directly to $h(\ell(t))$ yields $2^{w(w-1)/2}\cdot \mathrm{poly}(w)$, which is exponential in $w$, because the induction peels the $w$ chain elements one by one, accumulating a spurious doubling at each step.
% The curve $\Gamma$ is not an arbitrary curve in $(0,1)^w$: since $\sigma$ is strictly monotone and each $s_k$ depends on a single parameter $\lambda_k t + \mu_k$, each coordinate $s_k(t)$ is a monotone function of $t$, so $\Gamma$ is coordinate-wise monotone. Conceptually, this monotonicity is what allows the dimension-climbing argument in Proposition~\ref{prop:single-layer-degree} to bypass the per-coordinate Rolle step: the contribution of each chain element is read off directly from the behaviour of the frequency sum along the line (Lemma~\ref{lem:line-intersection}), rather than peeled off one at a time.
% \end{remark}

\begin{proposition}[Weaker single-layer bound]\label{prop:single-layer-mt}
Under the hypotheses of Proposition~\ref{prop:single-layer-degree}, one has
\begin{equation*}
  \md(V_{ij}) \;\leq\; C_{\mathrm{MT}}(n,L)\cdot w^{n^2},
\end{equation*}
where $C_{\mathrm{MT}}(n,L)$ depends only on $n$ and $L$.
\end{proposition}

\begin{proof}
By Remark~\ref{re:affine}, we may assume the regular value of the Gauss map is $v=e_n$, and bound the number of solutions of the system
\begin{equation}\label{eq:mt-gauss-system}
  f(x) = 0, \qquad \frac{\partial f}{\partial x_j}(x) = 0, \quad j = 1,\ldots, n-1.
\end{equation}
The pairwise difference $f = F_j - F_i$ has the form
\begin{equation}\label{eq:mt-f}
  f(x) = c_0 + \sum_{k=1}^w d_k\, \sigma(a_k\cdot x + b_k),
\end{equation}
where $a_k\in \R^n$ are the rows of $A^{(1)}$, $b_k\in \R$ are bias terms, and $d_k\in \R$ are output-layer weights. The partial derivatives are
\begin{equation*}
  \frac{\partial f}{\partial x_j}(x) = \sum_{k=1}^w d_k\, a_{kj}\, \sigma'(a_k\cdot x + b_k).
\end{equation*}
Each of the $n$ equations in~\eqref{eq:mt-gauss-system} has the form $c + \sum_k \gamma_k \sigma'(a_k\cdot x + b_k)$ (with the leading equation $f=0$ replacing $\sigma'$ by $\sigma$, which has the same Pfaffian form). By Lemma~\ref{lem:line-intersection}, each equation has at most $D = O(w^n)$ zeros along a generic line.

We bound $|Z_1\cap \cdots \cap Z_n|$, where $Z_i = \{h_i = 0\}$ are the $n$ hypersurfaces defined by~\eqref{eq:mt-gauss-system} and each $Z_i$ meets a generic line in at most $D$ points. More precisely, we prove by induction on $k\in \{1,\ldots,n\}$ the following claim: any intersection $Z_{i_1}\cap \cdots \cap Z_{i_k}$ of $k$ of the hypersurfaces, transverse to a generic affine subspace of dimension $k$, has at most $D^k$ points on that subspace. Applying this to the fixed set $\{Z_1,\ldots,Z_n\}$ with $k=n$ yields the stated bound. For the base case $k = 1$, a single hypersurface $Z_i$ meets a generic line in at most $D$ points. For the inductive step, assume the claim holds for $k-1$. Write $\Gamma := Z_{i_2}\cap \cdots \cap Z_{i_k}$ and assume WLOG (by a generic linear change of coordinates) that the ambient $k$-dimensional affine subspace is $\R^k$ and that $x_k$ is a Morse function on the smooth one-dimensional manifold $\Gamma$. For notational simplicity we describe the inductive step in the case $k=n$; the general case is identical.

For $t\in \R$, define $L_t = \{x_n = t\}$ and consider the function $\varphi(t) = |Z_2\cap \cdots \cap Z_n \cap L_t|$. Since $L_t\cong \R^{n-1}$ and the restrictions $h_2|_{L_t},\ldots,h_n|_{L_t}$ define $n - 1$ hypersurfaces in $L_t$, each meeting a generic line in at most $D$ points, the inductive hypothesis gives $\varphi(t) \leq D^{n-1}$ for all regular values of $t$.

The function $\varphi$ changes only at critical values of $x_n|_\Gamma$. At a non-degenerate critical point $p\in \Gamma$, either a new pair of intersection points with $L_t$ is created (a local minimum of $x_n|_\Gamma$) or an existing pair is destroyed (a local maximum); each critical point therefore changes $\varphi$ by exactly $\pm 2$. The critical points of $x_n|_\Gamma$ are the points of $\Gamma$ at which $\partial h_2/\partial x_n = \cdots = \partial h_n/\partial x_n = 0$, which together with $h_2 = \cdots = h_n = 0$ defines the critical locus. Each such condition adds one equation with the same line-intersection structure, so the number of critical points is at most $D^{n-1}$ by the same inductive bound.

Now, $Z_1\cap \Gamma$ is a finite set (for generic $h_1$). Each point of $Z_1 \cap \Gamma$ lies in some fiber $L_t$, and there are at most $\varphi(t) \leq D^{n-1}$ branches of $\Gamma$ passing through $L_t$. Along each branch, $h_1$ restricts to a smooth function of $t$ whose zeros correspond to crossings of $Z_1$. Since $Z_1$ meets a generic line in at most $D$ points, the restriction of $h_1$ to each branch has at most $D$ zeros, giving $|Z_1\cap \Gamma| \leq D^{n-1}\cdot D = D^n$.

Applying this with $D = O(w^n)$ (from Lemma~\ref{lem:line-intersection}) gives at most $D^n = O(w^{n^2})$ solutions of~\eqref{eq:mt-gauss-system}. Multiplying by a factor of $2$ for the two orientations of the unit normal yields the claim.
\end{proof}

\begin{remark}[Where the slack comes from]
The argument above is tight at each inductive step but loose as a whole. At each level of the dimension climb we use the per-line bound $D = O(w^n)$ of Lemma~\ref{lem:line-intersection} as a black box, and the induction multiplies these bounds together to give $D^n = w^{n^2}$. The Milnor--Thom framework treats the $n$ hypersurfaces as if they were independent, in the sense that it exploits no shared structure between the equations beyond the common per-line bound. What it misses is that all $n$ equations of~\eqref{eq:mt-gauss-system} are exponential polynomials \emph{in the same $w$ frequencies $\{\lambda_k\}_{k=1}^w$}: after clearing denominators, they become Laurent polynomials in $y = \Psi(x)$ that share a common Newton polytope $\conv(A)$ (the zonotope of Remark~\ref{rem:zonotope}). The mixed-volume bound of Bernstein's theorem captures this sharing and yields the sharp $O(w^n)$ bound of Proposition~\ref{prop:single-layer-degree}; the dimension climb, by contrast, has no mechanism for detecting a shared support and pays the full $w^{n(n-1)}$ factor for ignoring it.
\end{remark}
